% Figura extraída de "Tema 1.tex" (línea 363 del original)
\documentclass[12pt,border=1pt]{standalone}
\usepackage[spanish,es-tabla]{babel}
\usepackage[utf8]{inputenc}
\usepackage{amsmath}
\usepackage{xcolor}
\usepackage{tikz}
\usepackage{pgfplots}
\usepackage{circuitikz}
\usetikzlibrary{babel, circuits.ee}

\definecolor{applegreen}{rgb}{0.55, 0.71, 0.0}
\definecolor{airforceblue}{rgb}{0.36, 0.54, 0.66}
\definecolor{amber}{rgb}{1.0, 0.49, 0.0}

\newcommand{\diff}[3][1]{%
	\IfBlankTF{#1}{\frac{\mathrm{d}{#2}}{\mathrm{d}{#3}}}%
	{\frac{\mathrm{d}^{#1} {#2}}{\mathrm{d}{#3}^{#1}}}%
}%
\newcommand{\pdiff}[3][1]{%
	\IfBlankTF{#1}{\frac{\partial {#2}}{\partial {#3}}}%
	{\frac{\partial^{#1} {#2}}{\partial {#3}^{#1}}}%
}%

\begin{document}
\begin{tikzpicture}[scale=0.7]
    \begin{axis}[
        xlabel=$E$,
        ylabel=$f(E)$,
        domain=0:2,
        samples=100,
        axis lines=middle,
        ymin=0,
        ymax=1.1,
        xlabel style={
            at={(ticklabel* cs:1)},
            anchor=north west,
            xshift=0mm, % Ajusta esta distancia según tus necesidades
        },
        ylabel style={
            at={(ticklabel* cs:1.05)},
            anchor=north west,
            xshift=0mm, % Ajusta esta distancia según tus necesidades
        },
    ]
    \addplot[blue, ultra thick, smooth] {1/(1+e^((x-1)/(0.026)))};
     \draw[dashed, red] (axis cs:1,0) -- (axis cs:1,1);
     \draw [dashed, red] (axis cs:0,1) -- (axis cs:1,1);
    \node [label={$E_\mathrm{F}$} ] at (axis cs:0.91,0) {};
    \end{axis}
\end{tikzpicture}
\end{document}
